\section{Discussion}
\label{sec:discussion}

\paragraph{Paradoxe de la latence INT8, chiffré.} La quantification INT8 divise par~4 l'empreinte mémoire
des poids mais \emph{n'accélère pas} l'inférence sur le Cortex-M4 --- elle la ralentit. L'instrumentation
cycle-à-cycle du kernel (Section~\ref{sec:system}) permet de dire où : sur les deux jeux, le MAC entier
coûte en moyenne \num{6785}~cycles et la requantification \num{3777.25}, la déquantification n'en pesant que
\num{200}. Le total INT8 s'établit à \SI{74}{\micro\second} contre \SIrange{48}{50}{\micro\second} en FP32.
Deux causes se cumulent donc : le MAC entier n'a rien à gagner face à un FPU matériel déjà efficace en
simple précision, et la requantification --- un poste qui n'existe pas dans le chemin flottant --- consomme à
elle seule un tiers du budget. La piste directe reste le portage sur noyaux SIMD entiers (CMSIS-NN), qui
attaquerait le premier poste ; le second appelle plutôt une chaîne à échelles en puissance de deux.

\paragraph{Un paradoxe d'énergie s'ajoute au paradoxe de latence.} On pourrait espérer que l'INT8, à défaut
d'être plus rapide, soit plus sobre. La mesure dit le contraire : \SI{67.47}{\micro\joule} par inférence en
INT8 contre \SI{67.65}{\micro\joule} en FP32, soit un écart de \SI{0.18}{\micro\joule} pour une incertitude
combinée de \SI{2.34}{\micro\joule}. Il n'y a donc pas d'effet mesurable. La conséquence pratique est nette :
sur cette cible, \textbf{l'INT8 se justifie par la mémoire, ni par la vitesse ni par l'énergie}. Le levier
énergétique se trouve ailleurs --- dans la fréquence, où la marge de latence dont on dispose largement
s'échange contre de l'autonomie.

\paragraph{Honnêteté « mesuré » vs « émulé ».} Nous séparons rigoureusement deux natures de résultats. Les
parités FP32, l'effondrement INT8 legacy, la récupération par kernel v2, les latences, la RAM et l'énergie
sont \emph{mesurés sur carte réelle} (Sprints~36, 40, 46 à 50 et 53). Le diagnostic causal, l'échelle
d'ablation et le balayage de profondeur sont \emph{émulés bit-à-bit sur PC} (Sprints~39 et~47) : l'émulateur
reproduit exactement l'arithmétique du kernel C, ce qui en fait une référence fiable, mais ce n'est pas une
mesure carte. Depuis la campagne v2, la récupération n'appartient plus à la seconde catégorie : elle est
établie sur les quatre cellules per-canal, avec une parité gelée de \num{1.000} contre l'émulateur. Ce qui
reste non mesuré --- le format \texttt{q15} sur carte, l'énergie par mise à jour continue, le profil par
phase --- est reporté « à mesurer » plutôt qu'estimé.

\paragraph{Limite : l'oubli n'est pas éprouvé sur ces deux jeux.} Nos scénarios continus y sont courts
(trois tâches) et l'oubli moyen y reste inférieur à \num{0.01} : ces jeux ne mettent pas la régularisation
EWC à l'épreuve. L'oubli catastrophique a été mesuré ailleurs dans le projet, sur CWRU, où il atteint
\num{0.8475} pour EWC contre \num{0.8578} en apprentissage naïf. Comme ce jeu ne recouvre pas les deux
nôtres, nous nous gardons de mélanger les deux constats : la contribution de cet article porte sur le
\emph{portage} et la \emph{quantification}, pas sur l'efficacité de la régularisation contre l'oubli.

\paragraph{Portée.} Le résultat central --- une PTQ naïve peut détruire une tête de décision binaire, et une
calibration par-tenseur/par-canal (ou un format Q15) la restaure --- dépasse le cas EWC : il concerne toute
tête légère quantifiée en aval d'un extracteur, régime fréquent en TinyML.
